\newpage\handout[true]
{\textsc{Math 2106-B: Foundations of Mathematical Proof}}{Fall 2025}
{Homework 4}
{\textit{Prof.} \href{https://ceheitsch.github.io/webpage/}{Dr. Christine Heitsch}}{Due Date: September 26th (24 Hours Extended)}
{Math 2106-B: HW4}
{\large
  \noindent
  \textbf{Name:} \HandoutAuthor \smallskip \\

  \noindent
  \textbf{GTID:} 903743506 \smallskip \\

  \noindent
  %%% List anyone with whom you discussed the problems
  \textbf{Collaborators:} None. This material reflects independent preparation. \smallskip \\

  \noindent
  %%% List all resources used OTHER than the textbook, lecture notes, 
  %%% webwork problems, ICA questions, and previous homework assignments.
  \textbf{Outside resources:} This work was informed by MATH 2106 course materials 
  (ICAs: in-class activities and WebWork contents) from \HandoutInstitution\ (the course textbook is 
  Chartrand \parencite{chartrand_mathematical_2018}).   
  Outside references specifically include
  textbooks written by 
  Grimaldi \parencite{grimaldi_discrete_2004} (used in MATH 3012) and 
  Rosen \parencite{rosen_discrete_2019} (used in CS 2050). \smallskip \\

  \noindent
  \textbf{Notice}: The answers contain cross-references throughout this document as clickable hyperlinks in most PDF viewers. However, the Gradescope PDF view might not support clickable hyperlinks.
}\newpage
% \printbibliography[
  heading=bibintoc,   
  title={References}     % change “References” text if needed
]
% `heading=bibnumbered` if you prefer it numbered chapter (in \tableofcontents)
% `heading=bibintoc`    if you prefer it unnumbered but still in the ToC\newpage
% The following table presents the standard logical equivalences from these sources together with 
the author's own (or the student, myself, Jaehoon Song) clarifications and labels where applicable.


\begin{table}[h!]
  \centering
  \renewcommand{\arraystretch}{1.2}
  \setlength{\tabcolsep}{8pt}
  \begin{tabular}{|p{5.0cm}|p{5.0cm}|p{5.2cm}|}
    \hline
    \rowcolor[HTML]{E0E0E0}
    \multicolumn{2}{|l|}{\textbf{Logical Equivalences}} & \textbf{Name} \\
    \hline
    $\,p \AND q \equiv q \AND p$ & $\,p \OR q \equiv q \OR p$ & Commutative Law \\
    \hline
    $\, (p \AND q) \AND r \equiv p \AND (q \AND r)$ & $\, (p \OR q) \OR r \equiv p \OR (q \OR r)$ & Associative Law \\
    \hline
    $\,p \AND (q \OR r) \equiv (p \AND q) \OR (p \AND r)$ & $\,p \OR (q \AND r) \equiv (p \OR q) \AND (p \OR r)$ & Distributive Law \\
    \hline
    $\,\overline{p \AND q} \equiv \sNOT{p} \OR \sNOT{q}$ & $\,\overline{p \OR q} \equiv \sNOT{p} \AND \sNOT{q}$ & DeMorgan's Law \\
    \hline
    \multicolumn{2}{|l|}{$\,\overline{\overline{p}} \equiv p$} & Double Negation law \\
    \hline
    $\,p \AND \boldsymbol{T} \equiv p$ & $\,p \OR \boldsymbol{F} \equiv p$ & Identity Law \\
    \hline
    $\,p \OR \boldsymbol{T} \equiv \boldsymbol{T}$ & $\,p \AND \boldsymbol{F} \equiv \boldsymbol{F}$ & Domination Law \\
    \hline
    $\,p \AND p \equiv p$ & $\,p \OR p \equiv p$ & Idempotent Law \\
    \hline
    \multicolumn{2}{|l|}{$\,p \THEN q \equiv \sNOT{p} \OR q$} & Implication Equivalence \\
    \hline
    \multicolumn{2}{|l|}{$\,\overline{p \THEN q} \equiv p \AND \sNOT{q}$} & Negation of Implication \\
    \hline
    \multicolumn{2}{|l|}{$\,p \THEN q \equiv \sNOT{q} \THEN \sNOT{p}$} & Contrapositive Equivalence \\
    \hline
    \multicolumn{2}{|l|}{$\,p \OR \sNOT{p} \equiv \boldsymbol{T}$} & Law of Tautology \\
    \hline
    \multicolumn{2}{|l|}{$\,p \AND \sNOT{p} \equiv \boldsymbol{F}$} & Law of Contradiction \\
    \hline
  \end{tabular}
\end{table}\newpage
% 
The following definitions are presented from standard mathematical sources (\textbf{Sets}).
Clarifications and labels have been added where applicable.

\begin{define}{Basic Set Concepts}\phantomsection\label{def:basic-set-concepts} % \hyperref[def:basic-set-concepts]{YOUR_CONTEXT}
  \begin{enumerate}
    \item \textbf{Universal Set:} A universal set $\setUNIV$ is a
    set that contains all objects under consideration in a particular context or problem.
    \item \textbf{Empty Set:} The empty set, denoted $\setNULL$,
    is the unique set that contains no elements.
    \item \textbf{Subset:} Let $A$ and $B$ be sets. $A$ is a
    subset of $B$ ($A \subseteq B$) if (and only if)
    for all $x \in A$, $x \in B$.
    \item \textbf{Set Equality:} Let $A$ and $B$ be sets. The sets $A$ and $B$ are
    equal ($A = B$) if $A \subseteq B$ and $B \subseteq A$.
    \item \textbf{Power Set:} Let $A$ be a set. The power set of $A$, denoted $\sPOWER{A}$ or $2^A$,
    is the set of all subsets of $A$. That is, $\mathcal{P}(A) = \{S\in\setUNIV : S \subseteq A\}$.
  \end{enumerate}
\end{define}
The following definitions extend the understanding of set operations
and properties.


\begin{define}{Set Operations}\phantomsection\label{def:set-operations} % \hyperref[def:set-operations]{YOUR_CONTEXT}
  \begin{enumerate}
    \item \textbf{Complementary Set:} Let $A$ be a subset of the universal set $\setUNIV$.
    The complement of $A$, denoted $A^c$ or $\overline{A}$, is the set of all elements
    in $\setUNIV$ that are not in $A$. That is, $A^c = \{x \in \setUNIV : x \notin A\}$.
    \item \textbf{Union:} Let $A$ and $B$ be sets. The union of $A$ and $B$,
    denoted $A \cup B$, is the set of all elements that are in $A$, or in $B$, or in both.
    That is, $A \cup B = \{x : x \in A \lor x \in B\}$.
    \item \textbf{Intersection:} Let $A$ and $B$ be sets. The intersection of $A$ and $B$,
    denoted $A \cap B$, is the set of all elements that are in both $A$ and $B$.
    That is, $A \cap B = \{x : x \in A \land x \in B\}$.
    \item \textbf{Set Difference:} Let $A$, $B$ be subsets of a universal set $\setUNIV$.
    The set difference $A \sMINUS B$ is defined to be the set
    $\{ x \in \setUNIV \LDIV x \in A \text{ and } x \notin B\}$. \vspace*{10px}\\
    \textit{It is sometimes called the relative complement of $B$ in $A$,
    and denoted in the textbook as $A - B$.}
    \item \textbf{Cartesian Product:} Let $A$ and $B$ be sets. The Cartesian product of
    $A$ and $B$, denoted $A \sTIMES B$, is the set consisting of all ordered pairs:
    \[
    A \sTIMES B = \{(a, b) \LDIV a \in A, b \in B\}.
    \]
  \end{enumerate}
\end{define}


\newpage

The following claims are given in the reference as well. Let $A,B\in U$, a universal set.


\begin{claim}{Subset of Union: $A \subseteq \sOR{A}{B}$}\phantomsection\label{claim:subset-union} % \hyperref[claim:subset-union]{YOUR_CONTEXT}
  \begin{subprove}[bluecolor]
      Let $x\in U$, assume $x\in A$. Then the statement \inlinequote{$x\in A$ or $x\in B$} is true.
      Hence, by the definition, $\sOR{A}{B}=\{x\in U\LDIV x\in A$ or $x\in B\}$, $x\in \sOR{A}{B}$.
      Thus, $A \subseteq \sOR{A}{B}$.
  \end{subprove}
\end{claim}

\begin{claim}{Empty Subset: $\setNULL \subseteq A$}\phantomsection\label{claim:empty-subset} % \hyperref[claim:empty-subset]{YOUR_CONTEXT}
  \begin{subprove}[bluecolor]
      Let $x\in U$, assume $x\in \setNULL$. But there are no elements in $\setNULL$, so
      $x\in \setNULL$ is false. Since the premise is false, the implication if $x\in \setNULL$,
      then $x\in A$ is true. Thus, $\setNULL \subseteq A$ by definition of \hyperref[def:basic-set-concepts]{subsets}.
  \end{subprove}
\end{claim}

\begin{claim}{Subset as Equality: $\sOR{A}{B}=A$ if and only if $B\subseteq A$}\phantomsection\label{claim:subset-equality} % \hyperref[claim:subset-equality]{YOUR_CONTEXT}
  \begin{subprove}[bluecolor]
      Let $x\in U$, first, we assume $B\subseteq A$, we have already shown
      $A \subseteq \sOR{A}{B}$ by the claim, \hyperref[claim:subset-union]{Subset of Union}. 
      To show $\sOR{A}{B} \subseteq A$, assume $x\in \sOR{A}{B}$.
      By definition, $x\in A$ or $x\in B$. Considering $x\in B$ and assumption
      $B \subseteq A$, $x\in B$ implies $x\in A$. Thus, $\sOR{A}{B} \subseteq A$.
      Hence, if $B\subseteq A$, then $\sOR{A}{B} = A$. \vspace*{10px}\\
      Now, we assume $\sOR{A}{B}=A$. To show $B\subseteq A$, we now assume $x\in B$. Since $x\in B$, by definition
      of union, $x\in \sOR{A}{B} = A$. Therefore, $B\subseteq A$.
  \end{subprove}
\end{claim}


%%%%%%%%%%%%%%%%%%% 9/8/25

Let $\setUNIV$ be a set. 
\begin{claim}{ICA6}\phantomsection\label{claim:ica6} % \hyperref[claim:ica6]{YOUR_CONTEXT}
  For all $A,B\in\setUNIV$, $$A\sTIMES B =\setNULL\THEN \group{A=\setNULL\OR B=\setNULL}$$
  \begin{subprove}[bluecolor]
    Let $A,B\in\setUNIV$. We prove $A\sTIMES B=\setNULL$. Assume it is not the case that 
    $A=\setNULL$ or $B=\setNULL$. Thus, we have $A=\setNULL$ and $B=\setNULL$. 
    Since $A\ne\setNULL$ by assumption, there exists an element $a\in A$. Likewise, there 
    is some $b\in B$ since $B\ne\setNULL$. Thus, $\group{a,b}\in\sBUILD{(x,y)}{x\in A,y\in B}$ 
    by definition. But then, $A\sTIMES B\ne\setNULL$ by definition. Hence, the claim holds. 
  \end{subprove}
\end{claim}


\begin{define}{Relatively Prime Integers (Coprime)}\phantomsection\label{def:relatively-prime} % \hyperref[def:relatively-prime]{YOUR_CONTEXT}
  Let $x,y \in \mathbb{Z}$. The integers $x$ and $y$ are said to be
  \textbf{relatively prime} if and only if
  $\gcd(x,y) = 1$.
\end{define}
\newpage
% The following definitions and axioms are presented from standard mathematical sources (\textbf{Number Theory}) together with the author's own (or the student, myself, Jaehoon Song) clarifications and labels where applicable.

\begin{define}{Integers}[def:integers][red]
  $\setZ$: the set of all integers
\end{define}

\begin{define}{Rational Numbers}[def:rational-numbers][red]
  $\setQ$: the set of all rational numbers
\end{define}

\begin{define}{Real Numbers}[def:real-numbers][red]
  $\setR$: the set of all real numbers
\end{define}

\begin{define}{Natural Numbers}[def:natural-numbers][red]
  $\setN$ or $\setZ_{>0}$: the set of all natural numbers which are the set of positive integers (not including $0$)
\end{define}

\begin{define}{Positive Real Numbers}[def:positive-real-numbers][red]
  $\setR_{>0}$: the set of all positive real numbers
\end{define}

\begin{define}{Non-negative Real Numbers}[def:non-negative-real-numbers][red]
  $\setR_{\geq 0}$: the set of all non-negative real numbers
\end{define}


\newpage

The following definitions and claims extend the understanding of integers.

Here are the ICA2 Definitions and Axioms as follows.

\begin{define}{Even Integer}[def:even-integer][red]
  An integer $n$ is \underbar{even} if there exists 
  (existential quantification logic $\exists$) an integer $k$ such that
  \[
  n = 2k.
  \]
\end{define}

\begin{define}{Odd Integer}[def:odd-integer][red]
  An integer $n$ is \underbar{odd} if 
  there exists some integer $k$ such that
  \[
  n = 2k + 1.
  \]
\end{define}

\begin{define}{Closure under Addition}[def:closure-addition][red]
  The sum of two integers is again an 
  integer. (Closure under addition)
\end{define}

\begin{define}{Closure under Multiplication}[def:closure-multiplication][red]
  The product of two integers is again 
  an integer. (Closure under multiplication)
\end{define}

\begin{define}{Odd if and only if Not Even}[def:odd-iff-not-even][red]
  An integer is odd if and only if it 
  is not even.
\end{define}

Here are the ICA3 Definitions and Axioms as follows.

\begin{define}{Universal Closure under Addition}[def:universal-closure-addition][red]
  The sum of two integers (Universal quantification for all, for every, for each) is again an integer. (Closure under addition)
\end{define}

\begin{define}{Universal Closure under Multiplication}[def:universal-closure-multiplication][red]
  The product of two integers (Universal quantification for all, for every, for each) is again an integer. (Closure under multiplication)
\end{define}

\begin{define}{Even Integer Square}[def:even-integer-square][red]
  An integer $n$ is even if and only if $n^2$ is even.
\end{define}

\newpage

The following claims are given in the reference as well.

\begin{define}{Sum of Even Integers}[claim:sum-even-integers][red]
  The sum of two even integers is again an even integer.
  \begin{subprove}[red]
    Let $x$ be an integer and $y$ be an integer ($\forall x,y \in \setZ$). 
    Suppose $x$ is even and $y$ is even. Then, by \nameref{def:even-integer}, there exists an integer 
    $k$ and an integer $j$ such that $x=2k$ and $y=2j$. Then, 
    $x+y=\left(2k\right)+\left(2j\right)=2\left(k+j\right)$. By \nameref{def:closure-addition}, since $k$ and 
    $j$ are integers, then so is $\left(k+j\right)$. Hence, by definition, $x+y$ is an 
    even integer.
  \end{subprove}
\end{define}

The following definitions extend the understanding of integers including divisibility.

\begin{define}{Divisibility}[def:divisibility][red]
  An integer $n$ divides an integer $m$, denoted $n \LDIV m$, if there 
  exists an integer $k$ such that $m=n k$. Otherwise, $n \nmid m$.
\end{define}

The following definitions extend the classification of real numbers into rational and irrational numbers.

\begin{define}{Rational Number}[def:rational-number][red]
  A real number $x$ is \underbar{rational} if there exist integers $a$ and $b$ with $b \neq 0$ such that $bx = a$.
  Otherwise, $x$ is \underbar{irrational}.

  \textbf{Note}: Compare the rational number definition with the even/odd axiom \inlinequote{odd iff not even} from \nameref{def:odd-iff-not-even}.
\end{define}

The following claims are given in the reference as well.

\begin{define}{Irrationality of $\sqrt{2}$}[claim:sqrt2-irrational][red]
  $\sqrt{2}$ is irrational.
  \begin{subprove}[red]
    Suppose $\sqrt{2}$ is rational. Then, by \nameref{def:rational-number}, there exists $a,b\in \setZ$ 
    with $b\ne 0$ such that $b\cdot\sqrt{2}=a$. We may further assume that any factors 
    common to $a$ and $b$ have already been removed.\\

    Then, squaring both sides, this yields $2\cdot b^2 = a^2$.
    Since products of integers are again integers by \nameref{def:closure-multiplication}, we have that $a^2,b^2\in \setZ$.
    Thus, by definition, $a^2$ is even. Therefore, by \nameref{def:even-integer-square}, $a$ is even. 
    Thus, there exists $c\in\setZ$ such that $a=2c$. Then
    $$
    2b^2=a^2={2c}^2=4c^2
    $$
    Then, $b^2=2c^2$. Then, as before, $b^2$ is even, so $b$ is even.
    However, $a$ and $b$ both being even contradicts the fact that any factors 
    common to $a$ and $b$ have already been removed.
  \end{subprove}
\end{define}






%%%%%%%%%%%%%%%%%%% New style (refactor model)
\begin{define}{Relatively Prime Integers (Coprime)}[def:relatively-prime][red]
  Let $x,y \in \mathbb{Z}$. The integers $x$ and $y$ are said to be
  \textbf{relatively prime} if and only if
  $\gcd(x,y) = 1$.
\end{define}\newpage




%%%%%%%%%%%%%%%%%%%%%%%%%%%%%%%%%%%%%%%%%%%%%%%%%%%%%%%%%%%%%%%%%%%%%%%
% References
%%%%%%%%%%%%%%%%%%%%%%%%%%%%%%%%%%%%%%%%%%%%%%%%%%%%%%%%%%%%%%%%%%%%%%%
\textbf{Remark}: The following definitions (or results) are already proved in class or on previous ICA, 
Wbwk, Hmwk assignments, and contents of the textbook.





\begin{define}{Set Difference}[def:set-difference][red]
  Let $A$, $B$ be subsets of a universal set $\setUNIV$. The set difference $A \setminus B$ is defined to be the set $\{ x \in \setUNIV : x \in A \text{ and } x \notin B\}$.
\end{define}

\begin{define}{Pairwise Disjoint Sets}[def:pairwise-disjoint][red]
  Let $A_1, A_2, \ldots, A_n$ be sets. These sets are \textbf{pairwise disjoint} if for all $i \neq j$, we have $A_i \cap A_j = \emptyset$.
\end{define}

\begin{define}{Cartesian Product}[def:cartesian-product][red]
  Let $A$ and $B$ be sets. The Cartesian product of $A$ and $B$, denoted $A \times B$, is the set consisting of all ordered pairs:
  \[
  A \times B = \{(a, b) : a \in A, b \in B\}.
  \]
\end{define}






\begin{define}{Relation Properties}[def:relation-properties][red]
  Let $A$ be a set and $R$ a relation on $A$. Relations on a set are classified by their properties. The most important properties are:
  \begin{enumerate}
    \item \textbf{Reflexive}: The relation $R$ is reflexive if (and only if) for all $x \in A$, $(x, x) \in R$.
    \item \textbf{Symmetric}: The relation $R$ is symmetric if (and only if) for all $x, y \in A$, $(x, y) \in R$ implies $(y, x) \in R$.
    \item \textbf{Antisymmetric}: The relation $R$ is antisymmetric if (and only if) for all $x, y \in A$, $(x, y) \in R$ and $(y, x) \in R$ only if $x=y$.
    \item \textbf{Asymmetric}: The relation $R$ is asymmetric if (and only if) for all $x, y \in A$, if $(x, y) \in R$, then $(y, x) \notin R$.
    \item \textbf{Transitive}: The relation $R$ is transitive if (and only if) for all $x, y, z \in A$, if $(x, y) \in R$ and $(y, z) \in R$, then $(x, z) \in R$.
  \end{enumerate}
\end{define}

\begin{define}{Inverse Relation}[def:inverse-relation][red]
  Let $A$ and $B$ be sets. Let $R$ be a relation from $A$ to $B$. The inverse relation $R^{-1}$ from $B$ to $A$ is the set of ordered pairs $\{(b, a) \mid(a, b) \in R\}$.
\end{define}






% \begin{define}{Binary Relation}[def:binary-relation][red]
%   Let $A$ and $B$ be sets. A binary relation $R$ from $A$ to $B$ is a subset of $A \times B$. 
%   If (and only if) $$(a, b) \in R \subseteq A \times B$$, then we say 
%   that $a$ is related to $b$ and write $\relation{a}{b}$. Otherwise, $\nrelation{a}{b}$.
% \end{define}




% \begin{define}{Digraph}[def:digraph][red]
%   A digraph (or directed graph) is a set of vertices $V$ together with a set of edges $E \subseteq V \times V$.
% \end{define}

% \begin{define}{Bipartite Digraph}[def:bipartite-digraph][red]
%   A digraph is bipartite if (and only if) there exist subsets $A, B \subseteq V$ such that 
%   \begin{enumerate}
%     \item $V=A \cup B$ where $A \cap B=\emptyset$ and $A, B \neq \emptyset$.
%     \item for every $e \in E, e=(a, b)$ for some $a \in A, b \in B$.
%   \end{enumerate}
% \end{define}

% \textbf{Example}: Let $A=\{0,1,2\}$ and $B=\{a, b\}$. Then, $R=\{(0, a),(0, b),(1, a),(2, b)\}$ is a relation from $A$ to $B$. Binary relations on finite sets like this can be visualized as a bipartite digraph:

% \begin{figure}[h!]
%   \centering
%   \[
%   \xymatrix@R=1.5em@C=2.5em{ % Adjusted spacing to match the image
%     *++[o][F-]{0} \xyArc[r]\xyArc[ddr]  & *++[o][F-]{a} 
%     \\
%     *++[o][F-]{1} \xyArc[ur]            & 
%     \\
%     *++[o][F-]{2} \xyArc[r]             & *++[o][F-]{b}
%   }
%   \]
%   \caption{A bipartite graph example}
%   \label{fig:bipartite-graph}
% \end{figure}









% \begin{define}{Relation on a Set}[def:relation-on-set][red]
%   A relation on a set $\boldsymbol{A}$ is a relation from $A$ to $A$.
% \end{define}




% \begin{define}{Equivalence Classes Form Partition}[thm:equiv-classes-partition][red]
%   Let $R$ be an equivalence relation on a nonempty set $A$. Then, $P=\left\{[a]_R \mid a \in A\right\}$, the set of equivalence classes resulting from $R$, is a partition of $A$.
% \end{define}
% \begin{define}{Partition Induces Equivalence Relation}[thm:partition-equiv-relation][red]
%   Let $A$ be a nonempty set and $P=\left\{S_i \mid i \in I\right\}$ a partition of $A$. Then, there exists an equivalence relation $R$ which has the sets $S_i, i \in I$, as its equivalence classes.
% \end{define}

% \textbf{Remark}: Note that the set $I$ above is an arbitrary index set. It is not necessarily finite, countable, etc!


% \begin{define}{Modulo Operation}[def:modulo][red]
%   Let $a, m \in \mathbb{Z}$ with $m$ positive. The remainder when $a$ is divided by $m$ is denoted $\boldsymbol{a} \bmod \boldsymbol{m}$.
% \end{define}
% \begin{define}{Congruence Modulo}[def:congruence-modulo][red]
%   Let $a, b, m \in \mathbb{Z}$ with $m$ positive. If $m \mid(a-b)$ then we say that 
%   $\boldsymbol{a}$ is congruent to $\boldsymbol{b}$ modulo $\boldsymbol{m}$ and 
%   write $\boldsymbol{a} \equiv \boldsymbol{b}(\bmod \boldsymbol{m})$.
% \end{define}






% \textbf{Example}: Let $R$ be an equivalence relation on the set $\{1,2,3,4,5,6,7\}$ such that $1 R 3,3 R 4,4 R 7$, and $2 R 6$. Suppose that $R$ has three equivalence classes.



% \textbf{Example}: Let $S$ be another equivalence relation on the set $\{1,2,3,4,5,6\}$. Suppose the equivalence classes of $S$ are $\{1,4,5\},\{2,6\}$, and $\{3\}$.

% \textbf{Example}: Let $[a]_n$ denote the congruence class of $a \in \mathbb{Z}$ modulo $n$ (that is, the equivalence class of $a$ under the relation congruence modulo $n$ ) for integers $n \geq 2$. Are the following propositions true or false?







\begin{define}{Partition}[def:partition][red]
  Let $A$ be a nonempty set. A collection $\mathcal{P}$ of subsets of $A$ is defined as a \emph{partition} of $A$ if and only if the following conditions hold:
  \begin{itemize}
    \item[\textbf{1.}] \textbf{Non-emptiness:} No set in the collection is empty.
      $$ \forall S \in \mathcal{P}, \; S \neq \emptyset $$
    \item[\textbf{2.}] \textbf{Pairwise Disjointness:} The intersection of any two distinct sets in the collection is the empty set.
      $$ \forall S_1, S_2 \in \mathcal{P}, \; S_1 \neq S_2 \THEN S_1 \cap S_2 = \emptyset $$
    \item[\textbf{3.}] \textbf{Union:} The union of all sets in the collection constitutes the entire set $A$.
      $$ \bigcup_{S \in \mathcal{P}} S = A $$
  \end{itemize}
  The elements of $\mathcal{P}$ are referred to as the blocks or cells of the partition.
\end{define}

\begin{define}{Partition Refinement}[def:partition-refinement][red]
  Let $A$ be a nonempty set, and let $\mathcal{P}_1$ and $\mathcal{P}_2$ be 
  two partitions of $A$. The partition $\mathcal{P}_1$ is said to be a 
  \emph{refinement} of $\mathcal{P}_2$ if and only if every block of 
  $\mathcal{P}_1$ is a subset of some block of $\mathcal{P}_2$. 
  This relationship can be expressed symbolically as:
  $$ \forall S_1 \in \mathcal{P}_1, \; \exists S_2 \in \mathcal{P}_2 \text{ such that } S_1 \subseteq S_2 $$
\end{define}

\begin{define}{Equivalence Class}[def:equivalence-class][red]
  Let $R$ be an equivalence relation on a nonempty set $A$. For any element $a \in A$, 
  the \emph{equivalence class} of $a$ with respect to $R$, denoted $[a]_R$, is 
  defined as follows.
  $$ [a]_R = \{x \in A \mid (a,x) \in R \} $$
\end{define}








\newpage

%%%%%%%%%%%%%%%%%%%%%%%%%%%%%%%%%%%%%%%%%%%%%%%%%%%%%%%%%%%%%%%%%%%%%%%
% Problems
%%%%%%%%%%%%%%%%%%%%%%%%%%%%%%%%%%%%%%%%%%%%%%%%%%%%%%%%%%%%%%%%%%%%%%%
\begin{enumerate}
  \item Prove that if $R$ is an equivalence relation on a nonempty set $A$, then 
  the inverse relation $R^{-1}$ is also an equivalence relation on $A$.~\footnote{
    \textbf{Notice}: The answers contain cross-references throughout this document as clickable hyperlinks in most PDF viewers. 
    However, the Gradescope PDF view might not support clickable hyperlinks.
  }

  \begin{define}{Equivalence Relation}[def:equivalence-relation][red]
    Let $A$ be a nonempty set. A binary relation $R$ on $A$ (i.e., a subset $R \subseteq A \times A$) is designated an \emph{equivalence relation} if and only if it satisfies the following three axioms:
    \begin{itemize}
      \item[\textbf{1.}] \textbf{Reflexivity:} Every element in $A$ is related to itself.
        $$ \forall a \in A, \; (a,a) \in R $$
      \item[\textbf{2.}] \textbf{Symmetry:} If an element $a$ is related to an element $b$, then $b$ is also related to $a$.
        $$ \forall a,b \in A, \; (a,b) \in R \THEN (b,a) \in R $$
      \item[\textbf{3.}] \textbf{Transitivity:} If an element $a$ is related to $b$, and $b$ is related to $c$, then $a$ is related to $c$.
        $$ \forall a,b,c \in A, \; \left( (a,b) \in R \land (b,c) \in R \right) \THEN (a,c) \in R $$
    \end{itemize}
  \end{define}
  \begin{subprove}
    Let $R$ be an equivalence relation on a nonempty set $A$. 
    To show that $R^{-1}$ is an equivalence relation, it is required to verify that $R^{-1}$ satisfies reflexivity, symmetry, and transitivity by definition of \nameref{def:equivalence-relation}.
    \begin{itemize}
      \item \textbf{Reflexive:} Since $R$ is reflexive by definition, for all $a \in A$, $(a,a) \in R$. 
      Since $(a,a) \in R$, it follows that $(a,a) \in R^{-1}$ by definition of \nameref{def:inverse-relation}.
      \item \textbf{Symmetric:} Let $(a,b) \in R^{-1}$ be arbitrary. Then, $(b,a) \in R$ by definition. 
      Since $R$ is symmetric by definition, it follows that $(a,b) \in R$. 
      Therefore, $(b,a) \in R^{-1}$.
      \item \textbf{Transitive:} Let $(a,b), (b,c) \in R^{-1}$ be arbitrary. Then, $(b,a), (c,b) \in R$ by definition. 
      Since $R$ is transitive by definition, it follows that $(c,a) \in R$. 
      Therefore, $(a,c) \in R^{-1}$.
    \end{itemize}
    Therefore, $R^{-1}$ is an equivalence relation by definition of \nameref{def:equivalence-relation}.
  \end{subprove}





  

  \newpage

  \item Suppose that $R$ and $S$ are equivalence relations on the nonempty set $A$. Prove or disprove:~\footnote{
    \textbf{Notice}: The answers contain cross-references throughout this document as clickable hyperlinks in most PDF viewers. 
    However, the Gradescope PDF view might not support clickable hyperlinks.
  }
  \begin{enumerate}
    \item The set $R \sOR S$ is an equivalence relation.
      \begin{subdisprove}
        Let $A = \{1,2,3\}$ be a nonempty set. Let $R$ and $S$ be relations on $A$ where
        $R = \{(1,1),(2,2),(3,3),(1,2),(2,1)\}$ and
        $S = \{(2,2),(3,3),(2,3),(3,2)\}$. 
        Then, $(1,2) \in R \sOR S$ and $(2,3) \in R \sOR S$ by definition of Union, 
        but $(1,3) \notin R \sOR S$. 
        Since $(1,2), (2,3) \in R \sOR S$ but $(1,3) \notin R \sOR S$, 
        it follows that $R \sOR S$ is not transitive by definition of \nameref{def:relation-properties}.
        Therefore, $R \sOR S$ is not an equivalence relation by definition of \nameref{def:equivalence-relation}.
      \end{subdisprove}
    \item The set $R \sAND S$ is an equivalence relation.
      \begin{subprove}
        Let $R$ and $S$ be equivalence relations on a nonempty set $A$. 
        To show that $R \sAND S$ is an equivalence relation, it is required to verify that $R \sAND S$ satisfies reflexivity, symmetry, and transitivity by definition of \nameref{def:equivalence-relation}.
        \begin{itemize}
          \item \textbf{Reflexive:} Since $R$ and $S$ are reflexive by definition, 
          for all $a \in A$, $(a,a) \in R$ and $(a,a) \in S$. 
          Therefore, $(a,a) \in R \sAND S$ by definition of intersection.
          \item \textbf{Symmetric:} Let $(a,b) \in R \sAND S$ be arbitrary. Then, $(a,b) \in R$ and $(a,b) \in S$ by definition. 
          Since $R$ and $S$ are symmetric by definition, it follows that $(b,a) \in R$ and $(b,a) \in S$. 
          Therefore, $(b,a) \in R \sAND S$ by definition.
          \item \textbf{Transitive:} Let $(a,b), (b,c) \in R \sAND S$ be arbitrary. Then, $(a,b), (b,c) \in R$ and $(a,b), (b,c) \in S$ by definition. 
          Since $R$ and $S$ are transitive by definition, it follows that $(a,c) \in R$ and $(a,c) \in S$. 
          Therefore, $(a,c) \in R \sAND S$ by definition.
        \end{itemize}
        Therefore, $R \sAND S$ is an equivalence relation.
      \end{subprove}
  \end{enumerate}









  \newpage

  \item Let $R$ be a relation on a nonempty set $A$.~\footnote{
    \textbf{Notice}: The answers contain cross-references throughout this document as clickable hyperlinks in most PDF viewers. 
    However, the Gradescope PDF view might not support clickable hyperlinks.
  }
  \begin{enumerate}
    \item Prove that a relation $R$ on a set $A$ is symmetric if and only if $R=R^{-1}$.
      \begin{subprove}
        Let $R$ be a relation on a nonempty set $A$. 
        The proof is by showing both directions of the biconditional.
        \proofstep
        ($\Rightarrow$) Assume $R$ is symmetric. Let $(a,b) \in R$ be arbitrary. 
        Since $R$ is symmetric by definition of \nameref{def:relation-properties}, it follows that $(b,a) \in R$. 
        Therefore, $(a,b) \in R^{-1}$ by definition of \nameref{def:inverse-relation}. 
        Since $(a,b) \in R$ implies $(a,b) \in R^{-1}$, it follows that $R \subseteq R^{-1}$ by definition of subset.
        Conversely, let $(a,b) \in R^{-1}$ be arbitrary. Then, $(b,a) \in R$ by definition of \nameref{def:inverse-relation}. 
        Since $R$ is symmetric, it follows that $(a,b) \in R$. 
        Since $(a,b) \in R^{-1}$ implies $(a,b) \in R$, it follows that $R^{-1} \subseteq R$ by definition of subset.
        Therefore, $R = R^{-1}$.
        \proofstep
        ($\Leftarrow$) Assume $R = R^{-1}$. Let $(a,b) \in R$ be arbitrary. 
        Since $R = R^{-1}$, it follows that $(a,b) \in R^{-1}$. 
        Therefore, $(b,a) \in R$ by definition of \nameref{def:inverse-relation}. 
        Since $(a,b) \in R$ implies $(b,a) \in R$, it follows that $R$ is symmetric by definition.
        \proofstep
        Therefore, $R$ is symmetric if and only if $R = R^{-1}$.
      \end{subprove}
    \item Prove that a relation $R$ on a set $A$ is antisymmetric if and only if $R \sAND R^{-1} \subseteq \Delta$.
    \begin{define}{Diagonal Relation}[def:diagonal-relation][red]
      Let $A$ be a nonempty set. The \emph{diagonal relation} (or identity relation) on $A$, denoted $\Delta_A$ or $I_A$, is the relation of equality on the set $A$. It is formally defined as the set of all ordered pairs where both components are the same element from $A$:
      $$ \Delta_A = \{ (a,a) \mid a \in A \} $$
    \end{define}
      \begin{subprove}
        Let $R$ be a relation on a nonempty set $A$ and let $\Delta$ be the diagonal relation on $A$ by definition of \nameref{def:diagonal-relation}.
        The proof is by showing both directions of the biconditional.
        \proofstep
        ($\Rightarrow$) Assume $R$ is antisymmetric. Let $(a,b) \in R \sAND R^{-1}$ be arbitrary. 
        Then, $(a,b) \in R$ and $(a,b) \in R^{-1}$ by definition of intersection. 
        Since $(a,b) \in R^{-1}$, it follows that $(b,a) \in R$ by definition of \nameref{def:inverse-relation}. 
        Since $(a,b) \in R$ and $(b,a) \in R$, and $R$ is antisymmetric by definition of \nameref{def:relation-properties}, 
        it follows that $a = b$. Therefore, $(a,b) = (a,a) \in \Delta$ by definition of \nameref{def:diagonal-relation}.
        Since $(a,b) \in R \sAND R^{-1}$ implies $(a,b) \in \Delta$, it follows that $R \sAND R^{-1} \subseteq \Delta$ by definition of subset.
        \proofstep
        ($\Leftarrow$) Assume $R \sAND R^{-1} \subseteq \Delta$. Let $(a,b) \in R$ and $(b,a) \in R$ be arbitrary. 
        Since $(b,a) \in R$, it follows that $(a,b) \in R^{-1}$ by definition of \nameref{def:inverse-relation}. 
        Since $(a,b) \in R$ and $(a,b) \in R^{-1}$, it follows that $(a,b) \in R \sAND R^{-1}$ by definition. 
        Since $R \sAND R^{-1} \subseteq \Delta$, it follows that $(a,b) \in \Delta$ by definition of subset. 
        Therefore, $a = b$ by definition of \nameref{def:diagonal-relation}. 
        Since $(a,b) \in R$ and $(b,a) \in R$ implies $a = b$, it follows that $R$ is antisymmetric by definition.
        \proofstep
        Therefore, $R$ is antisymmetric if and only if $R \sAND R^{-1} \subseteq \Delta$.
      \end{subprove}
  \end{enumerate}


  \newpage
  \item Let $R$ be a relation on the nonempty set $A$. Prove or disprove:~\footnote{
    \textbf{Notice}: The answers contain cross-references throughout this document as clickable hyperlinks in most PDF viewers. 
    However, the Gradescope PDF view might not support clickable hyperlinks.
  }
  \begin{define}{Irreflexive Relation}[def:irreflexive-relation][red]
    Let $A$ be a nonempty set. A binary relation $R$ on $A$ is classified as \emph{irreflexive} if and only if no element in $A$ is related to itself. This property is stated formally as:
    $$ \forall a \in A, \; (a,a) \notin R $$
    Equivalently, a relation $R$ is irreflexive if its intersection with the diagonal relation $\Delta_A$ on $A$ is the empty set, i.e., $R \cap \Delta_A = \emptyset$.
  \end{define}
  \begin{enumerate}
    \item The relation $R$ could be neither reflexive nor irreflexive.
      \begin{subprove}
        Let $A = \{1,2\}$ be a nonempty set. Let $R = \{(1,1)\}$ be a relation on $A$. 
        Since $(1,1) \in R$ but $(2,2) \notin R$, it follows that $R$ is not reflexive by definition of \nameref{def:relation-properties}. 
        Since $(1,1) \in R$, it follows that $R$ is not irreflexive by definition of \nameref{def:irreflexive-relation}. 
        Therefore, $R$ is neither reflexive nor irreflexive.
      \end{subprove}
    \item The relation $R$ is reflexive if and only if $\sNOT{R}$ is irreflexive. 
      \begin{subprove}
        Let $R$ be a relation on a nonempty set $A$. 
        The proof is by showing both directions of the biconditional.
        \begin{itemize}
          \item ($\Rightarrow$) Assume $R$ is reflexive. Let $a \in A$ be arbitrary. 
          Since $R$ is reflexive by definition, it follows that $(a,a) \in R$. 
          Therefore, $(a,a) \notin \sNOT{R}$ by definition of complementary set. 
          Since for all $a \in A$, $(a,a) \notin \sNOT{R}$, it follows that $\sNOT{R}$ is irreflexive by definition of \nameref{def:irreflexive-relation}.
          \item ($\Leftarrow$) Assume $\sNOT{R}$ is irreflexive. Let $a \in A$ be arbitrary. 
          Since $\sNOT{R}$ is irreflexive by definition of \nameref{def:irreflexive-relation}, it follows that $(a,a) \notin \sNOT{R}$. 
          Therefore, $(a,a) \in R$. 
          Since for all $a \in A$, $(a,a) \in R$, it follows that $R$ is reflexive by definition.
        \end{itemize}
        Therefore, $R$ is reflexive if and only if $\sNOT{R}$ is irreflexive.
      \end{subprove}
    \item If $R$ is irreflexive, then so is $R^2=R \circ R$.
      \begin{subdisprove}
        Let $A = \{1,2\}$ be a nonempty set. Let $R = \{(1,2),(2,1)\}$ be a relation on $A$. 
        Since $(1,1) \notin R$ and $(2,2) \notin R$, it follows that $R$ is irreflexive by definition of \nameref{def:irreflexive-relation}. 
        However, $R^2 = R \circ R = \{(1,1),(2,2)\}$ by definition of composition. 
        Since $(1,1) \in R^2$ and $(2,2) \in R^2$, it follows that $R^2$ is not irreflexive by definition of \nameref{def:irreflexive-relation}. 
        Therefore, the statement is false by counterexample.
      \end{subdisprove}
  \end{enumerate}

  \newpage
  \item Let $R$, $S$, and $T$ be relations on $\setZ$ defined as:~\footnote{
    \textbf{Notice}: The answers contain cross-references throughout this document as clickable hyperlinks in most PDF viewers. 
    However, the Gradescope PDF view might not support clickable hyperlinks.
  }
  \begin{itemize}
    \item $a R b$ if and only if $a \equiv b \pmod{2}$,
    \item $a S b$ if and only if $a \equiv b \pmod{3}$,
    \item $a T b$ if and only if $a \equiv b \pmod{6}$.
  \end{itemize}
  \begin{enumerate}
    \item List the equivalence classes for $R$, $S$, and $T$.
      \begin{subanswer}
        Let $R$, $S$, and $T$ be the relations on $\setZ$ as defined above. The equivalence classes of $R$ are:
          \begin{itemize}
            \item $[0]_R = \{2k \mid k \in \setZ\}$
            \item $[1]_R = \{2k+1 \mid k \in \setZ\}$
          \end{itemize}
          The equivalence classes of $S$ are:
          \begin{itemize}
            \item $[0]_S = \{3k \mid k \in \setZ\}$
            \item $[1]_S = \{3k+1 \mid k \in \setZ\}$
            \item $[2]_S = \{3k+2 \mid k \in \setZ\}$
          \end{itemize}
          The equivalence classes of $T$ are:
          \begin{itemize}
            \item $[0]_T = \{6k \mid k \in \setZ\}$
            \item $[1]_T = \{6k+1 \mid k \in \setZ\}$
            \item $[2]_T = \{6k+2 \mid k \in \setZ\}$
            \item $[3]_T = \{6k+3 \mid k \in \setZ\}$
            \item $[4]_T = \{6k+4 \mid k \in \setZ\}$
            \item $[5]_T = \{6k+5 \mid k \in \setZ\}$
          \end{itemize}
      \end{subanswer}
    \item Find the equivalence classes of $R \sAND S$.
      \begin{subanswer}
        Let $R$ and $S$ be the relations on $\setZ$ as defined above. 
        Since $R \sAND S$ is the relation where $a (R \sAND S) b$ if and only if $a \equiv b \pmod{2}$ and $a \equiv b \pmod{3}$, 
        it follows that $R \sAND S$ is congruence modulo $6$ by the Chinese Remainder Theorem. 
        Therefore, the equivalence classes of $R \sAND S$ are:
        \begin{itemize}
          \item $[0]_{R \sAND S} = \{6k \mid k \in \setZ\}$
          \item $[1]_{R \sAND S} = \{6k+1 \mid k \in \setZ\}$
          \item $[2]_{R \sAND S} = \{6k+2 \mid k \in \setZ\}$
          \item $[3]_{R \sAND S} = \{6k+3 \mid k \in \setZ\}$
          \item $[4]_{R \sAND S} = \{6k+4 \mid k \in \setZ\}$
          \item $[5]_{R \sAND S} = \{6k+5 \mid k \in \setZ\}$
        \end{itemize}
      \end{subanswer}
    \item Prove the partition resulting from $T$ is a refinement of the partition resulting from $S$.
      \begin{subprove}
        Let $T$ and $S$ be the relations on $\setZ$ as defined above. 
        To show that the partition resulting from $T$ is a refinement of the partition resulting from $S$, 
        it is required to show that every equivalence class of $T$ is contained in some equivalence class of $S$ by definition of \nameref{def:partition-refinement}.
        Let $[k]_T$ be an arbitrary equivalence class of $T$ where $0 \leq k < 6$. 
        Since $a \in [k]_T$ if and only if $a \equiv k \pmod{6}$ by definition of \nameref{def:equivalence-class}, 
        and since $a \equiv k \pmod{6}$ implies $a \equiv k \pmod{3}$, 
        it follows that $a \in [k \bmod 3]_S$ by definition of \nameref{def:equivalence-class}. 
        Therefore, $[k]_T \subseteq [k \bmod 3]_S$ by definition of subset. 
        Since this holds for all equivalence classes of $T$, it follows that the partition resulting from $T$ is a refinement of the partition resulting from $S$ by definition of \nameref{def:partition-refinement}.
      \end{subprove}
    \item What is the relationship between the partition resulting from $T$ and the partition resulting from $R$? Justify your answer.
      \begin{subanswer}
        Let $T$ and $R$ be the relations on $\setZ$ as defined above. 
        The partition resulting from $T$ is a refinement of the partition resulting from $R$. 
        This is because for each equivalence class $[k]_T$ of $T$ where $0 \leq k < 6$, 
        since $a \in [k]_T$ if and only if $a \equiv k \pmod{6}$ by definition of \nameref{def:equivalence-class}, 
        and since $a \equiv k \pmod{6}$ implies $a \equiv k \pmod{2}$, 
        it follows that $a \in [k \bmod 2]_R$ by definition of \nameref{def:equivalence-class}. 
        Therefore, $[k]_T \subseteq [k \bmod 2]_R$ by definition of subset, 
        and the partition resulting from $T$ is a refinement of the partition resulting from $R$ by definition of \nameref{def:partition-refinement}.
      \end{subanswer}
    \item What is the relationship between $S$ and $T$? Between $R$ and $T$? Justify your answers.
      \begin{subanswer}
        Let $R$, $S$, and $T$ be the relations on $\setZ$ as defined above. 
        Since congruence modulo $6$ implies congruence modulo $3$ and congruence modulo $2$, 
        it follows that $T \subseteq S$ and $T \subseteq R$ by definition of subset. 
        However, $S \not\subseteq T$ and $R \not\subseteq T$ since congruence modulo $3$ does not imply congruence modulo $6$, 
        and congruence modulo $2$ does not imply congruence modulo $6$.
      \end{subanswer}
    \item Is there any relationship between $R$ and $S$? Justify your answers.
      \begin{subanswer}
        Let $R$ and $S$ be the relations on $\setZ$ as defined above. 
        Since congruence modulo $2$ does not imply congruence modulo $3$ and congruence modulo $3$ does not imply congruence modulo $2$, 
        it follows that $R \not\subseteq S$ and $S \not\subseteq R$ by definition of subset. 
        Since $R \neq S$ (they have different equivalence classes), 
        it follows that $R$ and $S$ are independent equivalence relations with no subset relationship between them.
      \end{subanswer}
  \end{enumerate}
\end{enumerate}